\chapter{The solution in one page}\label{ch:solution}

Near listens to one spoken report and keeps a running answer to three questions: where is the
casualty, how bad is it, and how do we get them out. Figure~\ref{fig:concept} shows the loop.

\nanofig{concept}{The Near loop: describe what you see, get located, answer one question, get the read-back.}
{Horizontal four-step infographic, 16:9, clean flat vector, plain white background, navy text, cyan and amber accents.
Step 1 ``DESCRIBE'': a soldier icon with a radio and a speech bubble showing a castle icon with ``11 o'clock'',
a church icon with ``240°'' and a sun icon with ``1 o'clock''. Step 2 ``LOCATE'': a top-down street map with a
large translucent cyan ellipse and two smaller dashed circles labelled A and B. Step 3 ``ASK'': a speech bubble
from a laptop saying ``Is the church left or right of the keep?'' with a small bar chart labelled ``bits gained''.
Step 4 ``READ BACK'': the ellipse shrunk to a small circle on a street corner, a red pin labelled ``P1'', a timer
``TQ 72:14'' and the text ``within 60 m''. Thin arrows between steps, a curved arrow from step 3 back to step 2.
Use exactly the quoted labels. Colour palette (use only these): navy \#0D1110, cyan \#1AA7C4, amber \#E09A2B, red \#E5484D, green \#3FA877, grey \#8B9A94, plus white for text on dark shapes.}

\section{Three users}

\paragraph{The field user} is a buddy or a medic, sometimes the casualty himself. He holds push-to-talk
and speaks a report in his own words:
\begin{quote}\itshape
``This is Bravo 2. One casualty, blast, left leg bleeding, tourniquet on at 14:32, conscious and
breathing. We're in a tree line. I see the castle keep at my 11 o'clock, a church at bearing 240,
sun at my 1 o'clock. Over.''
\end{quote}
He hears back:
\begin{quote}\itshape
``Bravo 2, this is CCP. Copy. One urgent casualty, tourniquet 14:32. Stretcher needed. Two possible
areas, best one within 300 metres. Question: if you have a compass, what is the bearing to the
castle keep? Degrees or mils. Over.''
\end{quote}
One to three exchanges later the position is confirmed, and if the casualty can walk, the read-back
turns into directions to the collection point.

\paragraph{The coordinator} sits at the casualty collection point (CCP) with a laptop. The map shows
every casualty with its uncertainty, the threat zones, the routes. Side tiles show the MIST card,
the condition with countdowns, the impact graph, the kit and team to send, and the evidence behind
the position (Figure~\ref{fig:dashboard}).

\paragraph{The radio operator} is often the coordinator. In radio mode the laptop listens to the
radio, transcribes each transmission and proposes the reply. The operator approves, edits or
skips it, and the laptop speaks it into the radio.

\nanofig{dashboard}{The coordinator dashboard: evacuation order, casualty card, map with uncertainty and threat
zones, location evidence, sound analysis, impact graph and the team to send.}
{High-fidelity dark-mode UI mockup of a tactical web dashboard, 16:9, the dashboard window itself in near-black navy tones on a plain white background,
thin 1 px borders, monospace numbers. Top bar: ``NEAR · COORDINATOR'', a red ``P1'' badge, KPI boxes
``POSITION ±60 m'', ``TOURNIQUET 72:14'' (amber), ``MOBILITY stretcher'', ``WORSE IN 4 min'' (red), buttons
``Radio'', ``Simulate''. Left column: ``EVACUATION ORDER'' list of three rows with P1, P1, P2 badges and timers.
Centre: a dark map of a small town with a castle, a cyan route line to a ``CCP'' marker, a translucent red circle
labelled ``SMALL-ARMS FIRE'', a pink dashed circle labelled ``DRONE'', a cyan ellipse around a red pin. Right
column: ``CASUALTY'' card with M A R C H rows coloured red/green/grey, ``LOCATION'' panel with a top-down building
plan and echo rings. Bottom: a left-to-right graph of boxes ``Reported → Body region → Possibly impacted → MARCH
threat → Bring''. Crisp, realistic product screenshot style. Colour palette (use only these): navy \#0D1110, cyan \#1AA7C4, amber \#E09A2B, red \#E5484D, green \#3FA877, grey \#8B9A94, plus white for text on dark shapes.}

\section{Three ways to run}

\begin{table}[htbp]\centering\small
\caption{Deployment modes. Near detects what is available and degrades gracefully.}\label{tab:modes}
\begin{tabularx}{\linewidth}{lXXX}
\toprule
& Connected & Radio only (war default) & Radio + offline phone\\
\midrule
Field equipment & Phone with a browser & Tactical voice radio & Radio + phone in airplane mode\\
Link to the CCP & Wi-Fi, 4G or a tunnel (HTTPS) & Voice channel & Voice channel carrying data tones\\
Speech-to-text & Local Whisper (or Chrome online) & Local Whisper on the laptop & Local Whisper on the laptop\\
Sensors used & Mic (sound, echoes, noise), compass, gyroscope, steps, optional last GPS fix & What is said; sounds heard through the radio & Phone measures offline, sends $\sim$30 B packets\\
Position clues & All & Clock directions, compass bearings, sun, ground, pace count & All, except the raw audio\\
Guidance & Spoken + live pin & Spoken over the radio + dead reckoning by time and pace count & On the phone: compass arrow, steps, spoken turns\\
\bottomrule
\end{tabularx}
\end{table}

\nanofig{modes}{The three deployment modes side by side.}
{Three-column technical diagram, 16:9, flat vector on white, navy lines, cyan and amber highlights. Column
``CONNECTED'': a phone icon linked by a Wi-Fi symbol to a laptop icon. Column ``RADIO ONLY'': a handheld tactical
radio linked by zig-zag radio waves to a second radio, which is linked by an audio cable to a laptop; a crossed-out
satellite icon and a crossed-out cloud icon above. Column ``RADIO + OFFLINE PHONE'': a phone held next to a radio,
small sound-wave tones drawn between the phone speaker and the radio microphone, a label ``~2 s of tones,
30 bytes'', radio waves to the laptop. Under each laptop a small map with a pin. Titles exactly ``CONNECTED'',
``RADIO ONLY'', ``RADIO + OFFLINE PHONE''. Colour palette (use only these): navy \#0D1110, cyan \#1AA7C4, amber \#E09A2B, red \#E5484D, green \#3FA877, grey \#8B9A94, plus white for text on dark shapes.}

\section{What the jury sees in the demo}

\begin{enumerate}
  \item An empty map of Vincennes on the coordinator laptop, drawn from local data with no internet.
  \item Outside by the Château, GPS off, a teammate plays the casualty; the other speaks the report on the radio.
  \item The map shows a large ellipse, sometimes two areas. The laptop asks, on the radio, the question
        that splits them.
  \item The answer comes in. The ellipse shrinks onto one street. Read-back: ``Position within 60 metres,
        priority 1''.
  \item A second unit reports a drone. A pink zone appears; the walking casualty's route bends around it.
  \item The evacuation list reorders live; the tourniquet timer keeps running.
\end{enumerate}
